\section{Translation: ST deconvolution as a coarsening problem}
\label{sec:translation}

\subsection{The DGP}

Consider a tissue containing cells of $K$ types. For each cell $i$
of type $k(i)$ and each gene $j$, the true expression $Y_{ij}$ is
drawn from a cell-type-specific distribution $f_{j,k(i)}(\cdot;
\theta_{j,k(i)})$ with mean $\mu_{j,k(i)}$. A spatial-transcriptomics
experiment partitions the tissue into $S$ spots; spot $s$ contains
$N_s$ cells with type composition vector $\bm{p}_s = (p_{s,1}, \ldots,
p_{s,K})$ where $p_{s,k}$ is the fraction of cells of type $k$ in
spot $s$. The observed spot-level count is approximately
\begin{equation}
  X_{sj} \approx N_s \sum_{k=1}^K p_{s,k} \mu_{j,k} + \text{technical noise},
  \label{eq:trans-spot}
\end{equation}
plus dropout in the same form as scRNA-seq
\citep{towell2026scrnacoarsening}.

The deconvolution problem is: given $\{X_{sj}\}_{s=1\ldots S,
j=1\ldots p}$ and an external scRNA-seq reference providing estimates
of $\mu_{j,k}$, recover the per-spot composition $\bm{p}_s$ and (when
desired) refine the per-cell-type expression estimates.

\subsection{The translation}

\Cref{tab:translation} states the explicit correspondence to
masked-data inference.

\begin{table}[ht]
\centering
\small
\begin{tabular}{@{}p{0.43\linewidth} p{0.53\linewidth}@{}}
\toprule
\textbf{Masked-data framework} & \textbf{Spatial transcriptomics} \\
\midrule
Component $j \in \{1,\ldots,m\}$ & Cell type $k \in \{1,\ldots,K\}$ \\
Component lifetime $T_{ij}$ & Per-cell expression $Y_{ij}$ for cell of type $k(i)$ \\
Component distribution $f_j(\cdot;\theta_j)$ & Cell-type expression model with mean $\mu_{j,k}$ \\
Failure system $T_i = \min$ & Spot total $X_{sj} \approx N_s \sum_k p_{s,k}\mu_{j,k}$ \\
Failed cause $K_i$ & Cell-type identity (latent for each cell) \\
\textbf{Candidate set} $c_i$ & \textbf{Spot-level cell-type composition} $\bm{p}_s$ \\
Singleton candidate $|c_i|=1$ & Single-cell-resolution platform (MERFISH/seqFISH+) \\
Identifying functions & Marker genes (cell-type-specific high-expression panel) \\
C1 (cause in candidate set) & Reference atlas covers all cell types present in tissue (no out-of-atlas types) \\
\textbf{C2} (coarsening at random): conditional on the candidate set, the masking distribution does not depend on which element is the true cause &
  Concretely: spot assignment (cell-to-spot membership) is determined by tissue geometry and protocol, not by the per-cell-type expression $\mu_{j,k}$ being measured. Violations include cell-type-dependent fixation or permeabilization differences that change which cells get captured \\
C3 (parameter independence) & $\bm{p}_s$ does not depend on the per-cell-type expression $\mu_{j,k}$ \\
\bottomrule
\end{tabular}
\caption{Translation between the masked-data framework
\citep{towell2026masked} and ST deconvolution.}
\label{tab:translation}
\end{table}

The DGP is a strict generalization of the scRNA-seq case from
\citet{towell2026scrnacoarsening}. There, the candidate set per
gene was binary ($\1\{X = 0\}$). Here, the candidate set per
\emph{spot} is a probability distribution over $K$ cell types,
parametrized by the composition vector $\bm{p}_s$. The richer
structure makes the rank condition more demanding (\cref{sec:identifiability})
but also gives more information per observation.

\subsection{Worked example: RCTD as a C1--C2--C3 specialization}
\label{sec:rctd-worked}

To make the ``subsumes as special case'' claim concrete, we derive
RCTD \citep{cable2021robust} as a parametric specialization of the
framework. RCTD models the spot-level UMI count $X_{sj}$ for gene
$j$ as Poisson with mean $N_s \sum_k p_{s,k} \mu_{j,k}$ (possibly
inflated by a multiplicative platform-effect term, which we omit
here without loss of generality). The likelihood factorizes across
spots and genes:
\begin{equation}
  \mathcal L(\bm \mu, P) \;=\; \prod_{s, j} \mathrm{Poisson}\!\left(X_{sj}
    \,;\, N_s \textstyle\sum_k p_{s,k}\mu_{j,k}\right).
  \label{eq:rctd-likelihood}
\end{equation}
RCTD fits $\hat{\bm p}_s$ by maximizing
\eqref{eq:rctd-likelihood} over $\bm p_s \in \Delta^{K-1}$ with
$\bm\mu$ fixed at scRNA-seq-reference estimates $\hat\mu_{j,k}$, no
penalty or prior on $\bm p_s$ beyond the simplex constraint.

The mapping into the framework is line-by-line. C1 (cause in
candidate set) is the assumption that every cell type present in a
spot appears in the reference $\hat\mu_{j,k}$: an out-of-atlas cell
type would be silently absorbed into the in-atlas types, biasing
$\hat{\bm p}_s$. A weaker but still consequential form of C1
violation is atlas-tissue distributional mismatch, in which a
cell type appears in both but with $\hat\mu_{j,k}$ from the
reference systematically displaced from the tissue's actual
$\mu_{j,k}$ for that type (a known issue in cross-study atlas
reuse; see \citet{tasic2018shared}). C2 (CAR) is the assumption that spot assignment
does not depend on per-cell-type expression: which cells fall in
which Visium spot is determined by physical tissue placement on the
slide, not by the gene-expression values being measured. C3
(parameter independence) is the assumption that $\bm p_s$ does not
encode information about $\bm \mu$: the cell-type composition of a
spot is a property of tissue architecture, separable from the
expression profile of each cell type. Under all three, the Poisson
likelihood in \eqref{eq:rctd-likelihood} is the ignorable-coarsening
likelihood of the framework with the parametric choice
$f_{j,k}(x; \mu_{j,k}) = \mathrm{Poisson}(x; N_s p_{s,k}\mu_{j,k})$.
The rank condition (\cref{thm:rank}) and cell-total consistency
(\cref{thm:cell-total-st}) then apply: RCTD inherits
identifiability when the spot-by-cell-type composition matrix
$P$ has full column rank $K$, and its fitted spot means match
observed spot means at any interior MLE.

The bias bound (\cref{thm:bias-marker}) bites for RCTD when the
reference $\hat\mu_{j,k}$ is restricted to a marker-gene panel
$\mathcal M$: $\hat{\bm p}_s$ inherits the first-order Taylor bias
$J_{p,M}\cdot\mathrm{vec}(\Delta M)$ with the sensitivity matrix
evaluated at the marker-restricted Hessian of
\eqref{eq:rctd-likelihood}.

Cell2location and Tangram fit the framework with different
parametric and procedural choices: Cell2location replaces the flat
simplex constraint with a hierarchical prior on $\bm p_s$ (a
specific form of identifiability assumption), and Tangram aligns
scRNA-seq cells to spatial positions, which the framework reads as
constructing singleton candidate sets at single-cell resolution.
Each method's identifiability and bias behavior follows from
substituting its specific parametric and procedural choices into the
same C1--C2--C3 ledger. We do not derive each in detail here; the
RCTD derivation above is the template.

CARD adds a spatial-smoothness prior on $\bm p_s$ that couples
neighboring spots; this falls outside the unconstrained-optimum
hypothesis of \cref{thm:cell-total-st} and is best analyzed as a
regularized estimator within the framework. The discussion
(\cref{sec:discussion}) enumerates the further parametric variants
(SPOTlight, Stereoscope, SpatialDecon, DestVI, STdeconvolve) that
fit the same template.
